% GENERATED FILE. Edit canonical lessons, then run npm run generate:books.
% canonical-source-hash: fnv1a64:308fd89009ec7eab
% canonical-lessons: TE-C366-khajana, TE-C366-nidhi, TE-C366-anvesana, TE-C366-dora, TE-C366-akramincu

\chapter{Treasure}
\label{ch:te-treasure}
\hlchaptermodality{366}

\begin{chapteropening}
\textbf{By the end of this chapter:} \emph{I can say a treasury, a treasure, a quest, a lord and to occupy.}

Complete the last lesson of chapter 366: I can say a treasury, a treasure, a quest, a lord and to occupy.
\end{chapteropening}

\section[\texorpdfstring{\te{ఖజానా} (khajānā)}{ఖజానా (khajānā)}]{\te{ఖజానా} (khajānā) — a treasury}

\label{lesson:TE-C366-khajana}



\begin{quote}
Before the new one: say the Telugu for to order, then the Telugu for to banish.
\end{quote}

\subsection*{You'll want to know: \te{ఖజానా}}
\textbf{\te{ఖజానా}} — \emph{khajānā} — \textquotedblleft{}a treasury\textquotedblright{}.

\begin{quote}
\te{రాజు} \te{ఖజానా} \te{బంగారంతో} \te{నిండి} \te{ఉండేది}.
\end{quote}

\begin{quote}
\emph{rāju khajānā baṅgārantō niṇḍi uṇḍēdi.}
\end{quote}

\begin{quote}
The king's treasury used to be full of gold.
\end{quote}

Today also the public purse: \te{ప్రభుత్వ} \te{ఖజానా}.

\begin{cousinweb}
\textbf{\te{ఖజానా}} (\emph{khajānā}) is a loan from Arabic \emph{khizāna}, \textquotedblleft{}a storehouse, a treasury\textquotedblright{}.
\end{cousinweb}

\begin{grammarlens}[title={What you've built}]
One more word for telling what happened.
\end{grammarlens}

\subsection*{Your turn}
\begin{itemize}
\raggedright
  \item \emph{Say it:} \emph{khajānā}
  \item \emph{Say it:} \emph{khajānā}, once more
  \item \emph{Say it:} \emph{rāju khajānā baṅgārantō niṇḍi uṇḍēdi.} — The king's treasury used to be full of gold.
  \item \emph{Recall:} say the Telugu for to order, then the Telugu for to banish, then say \emph{khajānā} again
\end{itemize}

\begin{tcolorbox}[breakable,colback=teal!4,colframe=teal!35!black,title={Before you move on}]
What does \te{ఖజానా} mean? (\textquotedblleft{}A treasury\textquotedblright{}.) Say it once more.
\end{tcolorbox}
\section[\texorpdfstring{\te{నిధి} (nidhi)}{నిధి (nidhi)}]{\te{నిధి} (nidhi) — a treasure; a fund}

\label{lesson:TE-C366-nidhi}



\begin{quote}
Before the new one: say the Telugu for to banish, then the Telugu for a treasury.
\end{quote}

\subsection*{You'll want to know: \te{నిధి}}
\textbf{\te{నిధి}} — \emph{nidhi} — \textquotedblleft{}a treasure; a fund\textquotedblright{}.

\begin{quote}
\te{పాత} \te{కోటలో} \te{ఒక} \te{నిధి} \te{ఉందని} \te{పుకారు} \te{ఉంది}.
\end{quote}

\begin{quote}
\emph{pāta kōṭalō oka nidhi undani pukāru undi.}
\end{quote}

\begin{quote}
There is a rumour that a treasure lies in the old fort.
\end{quote}

Also a fund: \te{విద్యా} \te{నిధి}, an education fund.

\begin{cousinweb}
\textbf{\te{నిధి}} (\emph{nidhi}) is borrowed from Sanskrit \emph{nidhi}, \textquotedblleft{}a setting down; a hoard\textquotedblright{} (CDIAL 7207).
\end{cousinweb}

\begin{grammarlens}[title={What you've built}]
One more word for telling what happened.
\end{grammarlens}

\subsection*{Your turn}
\begin{itemize}
\raggedright
  \item \emph{Say it:} \emph{nidhi}
  \item \emph{Say it:} \emph{nidhi}, once more
  \item \emph{Say it:} \emph{pāta kōṭalō oka nidhi undani pukāru undi.} — There is a rumour that a treasure lies in the old fort.
  \item \emph{Recall:} say the Telugu for to banish, then the Telugu for a treasury, then say \emph{nidhi} again
\end{itemize}

\begin{tcolorbox}[breakable,colback=teal!4,colframe=teal!35!black,title={Before you move on}]
What does \te{నిధి} mean? (\textquotedblleft{}A treasure; a fund\textquotedblright{}.) Say it once more.
\end{tcolorbox}
\section[\texorpdfstring{\te{అన్వేషణ} (anvēṣaṇa)}{అన్వేషణ (anvēṣaṇa)}]{\te{అన్వేషణ} (anvēṣaṇa) — a quest; a search}

\label{lesson:TE-C366-anvesana}



\begin{quote}
Before the new one: say the Telugu for a treasury, then the Telugu for a treasure.
\end{quote}

\subsection*{You'll want to know: \te{అన్వేషణ}}
\textbf{\te{అన్వేషణ}} — \emph{anvēṣaṇa} — \textquotedblleft{}a quest; a search\textquotedblright{}.

\begin{quote}
\te{నిధి} \te{కోసం} \te{అతని} \te{అన్వేషణ} \te{ఏళ్ళ} \te{తరబడి} \te{సాగింది}.
\end{quote}

\begin{quote}
\emph{nidhi kōsaṁ atani anvēṣaṇa ēḷḷa tarabaḍi sāgindi.}
\end{quote}

\begin{quote}
His search for the treasure went on for years.
\end{quote}

\begin{cousinweb}
\textbf{\te{అన్వేషణ}} (\emph{anvēṣaṇa}) is borrowed from Sanskrit \emph{anveṣaṇa}, \textquotedblleft{}seeking for, searching, investigating\textquotedblright{} (Monier-Williams).
\end{cousinweb}

\begin{grammarlens}[title={What you've built}]
One more word for telling what happened.
\end{grammarlens}

\subsection*{Your turn}
\begin{itemize}
\raggedright
  \item \emph{Say it:} \emph{anvēṣaṇa}
  \item \emph{Say it:} \emph{anvēṣaṇa}, once more
  \item \emph{Say it:} \emph{nidhi kōsaṁ atani anvēṣaṇa ēḷḷa tarabaḍi sāgindi.} — His search for the treasure went on for years.
  \item \emph{Recall:} say the Telugu for a treasury, then the Telugu for a treasure, then say \emph{anvēṣaṇa} again
\end{itemize}

\begin{tcolorbox}[breakable,colback=teal!4,colframe=teal!35!black,title={Before you move on}]
What does \te{అన్వేషణ} mean? (\textquotedblleft{}A quest; a search\textquotedblright{}.) Say it once more.
\end{tcolorbox}
\section[\texorpdfstring{\te{దొర} (dora)}{దొర (dora)}]{\te{దొర} (dora) — a lord; a master}

\label{lesson:TE-C366-dora}



\begin{quote}
Before the new one: say the Telugu for a treasure, then the Telugu for a quest.
\end{quote}

\subsection*{You'll want to know: \te{దొర}}
\textbf{\te{దొర}} — \emph{dora} — \textquotedblleft{}a lord; a master\textquotedblright{}.

\begin{quote}
\te{పూర్వం} \te{దొర} \te{గ్రామం} \te{అంతా} \te{ఏలేవాడు}.
\end{quote}

\begin{quote}
\emph{pūrvaṁ dora grāmaṁ antā ēlēvāḍu.}
\end{quote}

\begin{quote}
In the old days the lord ruled the whole village.
\end{quote}

In the villages \te{దొర} was the landlord, and \te{దొరసాని} his wife.

\begin{cousinweb}
\textbf{\te{దొర}} (\emph{dora}) is a Telugu word for a lord or master; its origin is not settled.
\end{cousinweb}

\begin{grammarlens}[title={What you've built}]
One more word for telling what happened.
\end{grammarlens}

\subsection*{Your turn}
\begin{itemize}
\raggedright
  \item \emph{Say it:} \emph{dora}
  \item \emph{Say it:} \emph{dora}, once more
  \item \emph{Say it:} \emph{pūrvaṁ dora grāmaṁ antā ēlēvāḍu.} — In the old days the lord ruled the whole village.
  \item \emph{Recall:} say the Telugu for a treasure, then the Telugu for a quest, then say \emph{dora} again
\end{itemize}

\begin{tcolorbox}[breakable,colback=teal!4,colframe=teal!35!black,title={Before you move on}]
What does \te{దొర} mean? (\textquotedblleft{}A lord; a master\textquotedblright{}.) Say it once more.
\end{tcolorbox}
\section[\texorpdfstring{\te{ఆక్రమించు} (ākramiñcu)}{ఆక్రమించు (ākramiñcu)}]{\te{ఆక్రమించు} (ākramiñcu) — to occupy; to invade}

\label{lesson:TE-C366-akramincu}



\begin{quote}
Before the new one: say the Telugu for a quest, then the Telugu for a lord.
\end{quote}

\subsection*{You'll want to know: \te{ఆక్రమించు}}
\textbf{\te{ఆక్రమించు}} — \emph{ākramiñcu} — \textquotedblleft{}to occupy; to invade\textquotedblright{}.

\begin{quote}
\te{శత్రువులు} \te{కోటను} \te{ఆక్రమించారు}.
\end{quote}

\begin{quote}
\emph{śatruvulu kōṭanu ākramiñcāru.}
\end{quote}

\begin{quote}
The enemies occupied the fort.
\end{quote}

Past (I …): \te{ఆక్రమించాను}, \emph{ākramiñcānu}.

Also to take land unlawfully: \te{స్థలం} \te{ఆక్రమించు}.

\begin{cousinweb}
\textbf{\te{ఆక్రమించు}} (\emph{ākramiñcu}) is built on Sanskrit \emph{ākramaṇa}, \textquotedblleft{}stepping on, approaching\textquotedblright{} (CDIAL 1016), from \emph{kram}, \textquotedblleft{}to step\textquotedblright{}, with the Telugu verb ending \emph{-iñcu}.
\end{cousinweb}

\begin{grammarlens}[title={What you've built}]
One more word for telling what happened.
\end{grammarlens}

\subsection*{Your turn}
\begin{itemize}
\raggedright
  \item \emph{Say it:} \emph{ākramiñcu}
  \item \emph{Say it:} \emph{ākramiñcu}, once more
  \item \emph{Say it:} \emph{śatruvulu kōṭanu ākramiñcāru.} — The enemies occupied the fort.
  \item \emph{Recall:} say the Telugu for a quest, then the Telugu for a lord, then say \emph{ākramiñcu} again
\end{itemize}

\begin{tcolorbox}[breakable,colback=teal!4,colframe=teal!35!black,title={Before you move on}]
What does \te{ఆక్రమించు} mean? (\textquotedblleft{}To occupy; to invade\textquotedblright{}.) Say it once more.
\end{tcolorbox}
